\begin{frame}{확인 문제: MDP인가, 마르코프인가? (1/2)}
  \kq{네 시나리오 각각에 대해 (a) MDP로 모델링할 수 있는가,
  (b) 가능하다면 $\mathcal{S}, \mathcal{A}, P, R$ 을 한 줄씩 쓰라:}
  \vskip0.6em
  \begin{itemize}
    \item \textbf{(a) 5층 건물의 엘리베이터 호출 제어}
      \begin{itemize}
        \item 답: \kb{MDP} --- $\mathcal{S} =$ (각 층의 상/하 호출
          버튼 상태 $2^{10}$ 조합) $\times$ (엘리베이터 현재 층 5)
          $=$ 5120개 이하, $\mathcal{A} = \{up, down, open\}$,
          $P$ = 버튼 상태 + 층 이동, $R = -1$(매 스텝 벌금) 또는
          ``누락 호출'' 페널티.
        \item 마르코프 성질은 \textbf{성립}(현재 호출 상태가 미래의
          모든 것을 결정).
      \end{itemize}
    \item \textbf{(b) ``오늘 기분''을 몰라서, 어제 몇 시간 폰을 봤는
      것만 관측으로 가지는 추천 문제}
      \begin{itemize}
        \item 답: \kb{마르코프 성질이 깨진} 문제 --- ``기분''이
          \textbf{hidden state}.
        \item 어제 폰 사용 시간만으로는 같은 관측 아래
          ``피곤한 날''과 ``재밌는 날''을 구별 못 해 \textbf{추천의
          결과 분포가 달라진다}.
        \item 해결: 최근 며칠의 관측을 확장 상태로 묶거나
          (전략 1), 사용자 모델의 은닉 상태로 요약(전략 2).
      \end{itemize}
  \end{itemize}
\end{frame}
